\section{Conclusion}
\label{sec:conclusion}

\looseness=-1 We identify conditioning-induced value overestimation (CIVO) as a systematic failure mode of temporal conditioning in imagination-based actor-critic training. Trained $\dt$-embeddings with multi-rate data exposure reliably induce value overestimation while frozen embeddings never do ($p < 0.001$), revealing co-adaptation between the temporal embedding and the world model as the causal factor. The NVE diagnostic achieves perfect binary classification (AUROC${=}1.000$, $n{=}36$), and cross-environment analysis localizes onset to moderate-complexity locomotion (6~DoF), where dynamics are complex enough for embedding co-adaptation yet baselines remain stable enough for CIVO isolation.

\looseness=-1 A four-level intervention hierarchy reveals systematic resilience: gradient interventions fail from pathway redundancy, output interventions fail from anchor co-drift, representation interventions exhibit phase transitions, and dual critics converge in pure imagination ($\rho > 0.999$), so Double Q pessimism is structurally ineffective in any system with shared imagined rollouts. Only wide dual-critic (min-aggregation $+$ increased capacity) achieves partial prevention, and even this result does not transfer across environments. For practitioners, we recommend:
\begin{enumerate}[leftmargin=1.5em,itemsep=1pt,topsep=2pt]
\item Validate with frozen-embedding controls before deploying learned temporal embeddings.
\item Monitor NVE as an early warning diagnostic.
\item Prefer MPC-style re-planning for variable-rate execution, where real-data correction prevents bias compounding.
\end{enumerate} Code and configurations will be released upon acceptance.

\section{Limitations}
\label{sec:limitations}

\looseness=-1 All experiments use DreamerV3 as the sole world-model architecture; whether CIVO manifests in transformer-based world models \citep{micheli2023iris} or diffusion-based planners remains untested and would establish whether the phenomenon generalizes beyond RSSM architectures. Observations are proprioceptive in simulation; pixel inputs and sim-to-real transfer have not been evaluated. Our compute budget (${\sim}80$ GPU-hours on a single GPU) restricts the factorial design scale; larger sweeps may reveal additional interaction effects. Per-environment onset rates carry wide confidence intervals (Walker 3/5: Clopper-Pearson 95\% CI $[15\%, 95\%]$; Cheetah 3/4: $[19\%, 99\%]$); the trained-vs-control separation (5/7 vs.\ 0/16, Fisher exact $p < 0.001$) is robust, but per-environment rates require more seeds to narrow. The NVE $2.0\times$ threshold was selected on training data; leave-one-environment-out cross-validation partially mitigates overfitting (\S\ref{sec:appendix_sensitivity}), but independent calibration on new environments is needed. RSSM fidelity analysis (Table~\ref{tab:rssm_fidelity}) is single-seed single-checkpoint; multi-seed verification would strengthen the mechanistic claims. The DoF-dependent onset boundary (\S\ref{sec:discussion}) has been mapped only for DMC locomotion; generalization to other domains and physical robots remains untested.
